Nanobodies are compact antibody fragments that can recognize specific protein surfaces. Designing them against a chosen epitope could support targeted therapeutics and molecular detection. The central problem is to identify sequences that both retain the nanobody fold and form the intended binding interface.\\\\
Sequence generation alone does not explicitly establish a binding geometry. Conversely, a plausible backbone may lack a sequence that reliably adopts it. Existing structure-based methods address these problems, but candidate selection remains uncertain: prediction confidence measures structural confidence, not experimental binding success. \href{https://www.ebi.ac.uk/training/online/courses/alphafold/inputs-and-outputs/evaluating-alphafolds-predicted-structures-using-confidence-scores/confidence-scores-in-alphafold-multimer/}{AlphaFold confidence documentation}\\\\
We propose to implement and evaluate an epitope-directed design pipeline built on existing models. The intuition is to constrain the binding geometry first, sample compatible sequences, and then check whether predicted structures recover the intended interface. The research contribution will be a controlled evaluation of design and selection choices within this established approach.
